\documentclass[11pt,a4paper]{article}
\usepackage[T2A]{fontenc}
\usepackage[utf8]{inputenc}
\usepackage[russian,english]{babel}
\usepackage{amsmath,amssymb,amsthm,mathtools}
\usepackage{setspace}
\usepackage{enumitem}
\usepackage[hidelinks]{hyperref}
\setstretch{1.3}
\frenchspacing
\hyphenpenalty=10000
\renewcommand{\leq}{\leqslant}
\renewcommand{\le}{\leqslant}
\renewcommand{\geq}{\geqslant}
\renewcommand{\ge}{\geqslant}
\newcommand{\head}[1]{\par\vspace{15pt}\noindent\textbf{\underline{#1}}\quad\ignorespaces}
\newcommand{\src}[1]{\unskip\nobreak\hfil\penalty50\hskip1em\hbox{}\nobreak\hfill(#1){\parfillskip=0pt\par}}
\begin{document}
\begin{center}
{\huge\bfseries Methods of proof -- 3\par}\vspace{1.1em}
{\Large The extremal principle and compactness\par}\vspace{1.1em}
{\Large Homework\par}\vspace{1.1em}
{\foreignlanguage{russian}{13 сентября 2026 года}\par}
\end{center}\vspace{0.0pt}

\noindent\textbf{\underline{Theoretical minimum.}}\quad The extremal principle is one line long:
take the object that maximises something, and use that it cannot be improved. Choosing the
functional is the whole difficulty, and the three usual choices are a size, a count and a
distance. On infinite families the existence of the extremum has to be earned: through
compactness and continuity, through a convergent subsequence, through Arzel\`a--Ascoli for
families of functions, or through Zorn's lemma when no topology is available. The variational
method is the same idea used constructively: maximise a functional and read off the conditions at
the maximiser. If the supremum is not attained, either the family must be enlarged so that it
becomes compact, or the limit object must be shown to be admissible~--- and this check is
precisely the step that most wrong solutions omit.

\medskip
\noindent\rule{\textwidth}{0.4pt}

\begin{enumerate}[itemsep=0.4em,topsep=7.4pt]
\item Prove that among all triangles inscribed in a given convex body there is one of maximal
area, and that its vertices have the property that the tangent lines are parallel to the opposite
sides.

\item Prove that a finite set of $n\ge5$ points in the plane, no three collinear, has a
subset of four points forming a convex quadrilateral, by taking the convex hull and arguing
extremally (compare with Convex -- 1).

\item Let $f$ be continuous on $\mathbb R$ with $f(x)\to0$ as $|x|\to\infty$ and $f$ not
identically zero. Prove that $f$ attains either its maximum or its minimum, and give an example
where only one of the two is attained.

\item Prove that every bounded infinite subset of $\mathbb R^n$ has an accumulation point, and
deduce that a sequence of nested non-empty compact sets has a common point; then give an example
of nested closed unbounded sets with empty intersection.

\item Prove that in a finite metric space there exist two points at maximal distance, and that
this pair is not necessarily unique; then prove that a compact metric space has a pair realising
its diameter.

\item Prove the existence of a best uniform approximation to a continuous function by polynomials
of degree at most $n$, and of a best $L^2$ approximation, and prove that both are unique,
explaining why the two uniqueness proofs are different (compare with Continuous functions -- 9).

\item Prove that a non-constant polynomial attains its minimum modulus on $\mathbb C$, and deduce
the fundamental theorem of algebra by showing that the minimum cannot be non-zero (d'Alembert's
argument).

\item Let $\mathcal A$ be a family of open intervals covering $[0,1]$. Prove by the extremal
principle that a finite subfamily covers $[0,1]$, and compare this with the standard compactness
proof.

\item Prove that among all closed curves of a given length the circle encloses the largest area,
at least in the class of convex curves, by a compactness plus symmetrisation argument, and say
where the existence of the maximiser is used (Convex -- 4).

\item Let $T$ be an infinite tree in which every vertex has finite degree. Prove that $T$ has an
infinite path (K\H{o}nig's lemma), and deduce the compactness principle for propositional logic
in the countable case.
\end{enumerate}
\end{document}
